\section*{Chapter \ref{ch:iv:the-final-round-and-trading-games} --- The Final Round and Trading Games}

\begin{solution}{iq:iv:the-final-round-and-trading-games:1}
Centre 21: each card averages 7. Width: the sum's standard deviation is $\sqrt{3 \times 14 \times 49/51}
\approx 6.35$; a market of 19 at 23 (two points each side, a third of a standard deviation) is tight enough to
trade and wide enough to absorb an informed trade or two. Say that you will move it as cards appear or as
trades come in.
\iqlookfor{a centred price from the mean, a width tied to the spread, and a stated plan to update.}
\end{solution}

\begin{solution}{iq:iv:the-final-round-and-trading-games:2}
Tonight: confirm the schedule and the interview types, prepare the three one-minute answers (background, why
this role, one piece of work in depth), revise the core families for the role, not new topics, and sleep.
Tomorrow: arrive early, eat and drink between interviews, treat each interview as independent, note the
questions after each one, and keep two good questions for each interviewer. If one interview goes badly, say
nothing about it in the next room.
\iqlookfor{preparation that matches the day's structure and management of energy across it.}
\end{solution}

\begin{solution}{iq:iv:the-final-round-and-trading-games:3}
The king is one of the three; the other two are drawn from 51 cards whose values sum to $364 - 13 = 351$, mean
$351/51 \approx 6.88$. Fair value $13 + 2 \times 6.88 \approx 26.76$.
\iqlookfor{conditioning on the revealed card and removing it from the deck.}
\end{solution}

\begin{solution}{iq:iv:the-final-round-and-trading-games:4}
Given the rule, the interviewer's card is uniform on 8 to 13 (mean 10.5), and the other two cards average
$(364 - c)/51$ each. Averaging $c + 2(364 - c)/51$ over $c = 8, \dots, 13$ gives about 24.36. The buy was
informative: raise the market to around 23 at 27 or tighter around 24.4, and be wary of selling more at 23
to the same player.
\iqlookfor{a Bayesian update from the counterparty's trading rule, not from the trade price alone.}
\end{solution}

\begin{solution}{iq:iv:the-final-round-and-trading-games:5}
The Kelly fraction at even money is $2p - 1 = 0.2$: bet 20\% of current chips each time (One Quant Book~2,
chapter~29). The expected log growth per bet is $0.6\ln 1.2 + 0.4\ln 0.8 \approx 0.0201$, so after ten bets
the median outcome is about $100e^{0.201} \approx 122$ chips; betting 10\% or 30\% grows more slowly. The mean
final stack is higher than the median and larger bets raise it, at the price of ruinous paths.
\iqlookfor{the Kelly fraction, the log-growth rate and the difference between mean and median wealth.}
\end{solution}

\begin{solution}{iq:iv:the-final-round-and-trading-games:6}
Take the structuring role without taking over: ``Twenty years is about 5\,000 trading days. Daily volatility
is about 1\% in calm years and 2 to 3\% in crises; if returns were normal with 1.2\% daily volatility, a 2\% move
is a 1.7-sigma event in either direction, about 10\% of days, but fat tails and volatility clustering change that. Can we each
give a number and a reason in thirty seconds, then combine?'' Then ask the quietest member for theirs, write
the numbers down, and summarise the decision at the end. The assessors score the structure, the number and the
inclusion, not whether the group's final answer is right.
\iqlookfor{structure, a first quantitative anchor, inclusion of others and a clear close.}
\end{solution}

\begin{solution}{iq:iv:the-final-round-and-trading-games:7}
With 5 and 6 shown, the third card is one of the 50 remaining, whose values sum to $364 - 11 = 353$, mean
$7.06$: fair value $11 + 7.06 = 18.06$. The position is worth $3 \times (18.06 - 24) \approx -17.82$ in
expectation. A bid of 18 is 0.06 below fair: selling costs about 0.06 a contract in expectation and removes the
remaining variance (one card, standard deviation about 3.8). Sell if the position is large against your
limits or the game rewards risk management; otherwise holding is marginally better. Say both.
\iqlookfor{the conditional fair value, mark-to-market, and a risk-aware decision stated with its trade-off.}
\end{solution}

\begin{solution}{iq:iv:the-final-round-and-trading-games:8}
The informed player's fair value is $F(c) = c + 2(364 - c)/51$, between $15.24$ (an ace) and $26.76$ (a king).
With a market of half-width $h$ around 21 the player buys when $F(c) > 21 + h$ and sells when $F(c) < 21 -
h$, and your expected loss is the average of the amounts by which $F$ crosses your quotes: about 3.10 per
opportunity at $h = 0$ and 1.43 at $h = 2$ (exact enumeration over the 13 card values). The player never
trades once $h \ge 98/17 \approx 5.76$. The cost of informed flow falls quickly with width; the game is about
choosing a width at which uninformed trades pay for the informed ones.
\iqlookfor{the informed player's conditional values, the loss as a function of width, and the break-even width.}
\end{solution}

\begin{solution}{iq:iv:the-final-round-and-trading-games:9}
Every rule wins with probability $30/52 = 15/26$. The proportion of red cards among those left is a
martingale as the deck is turned over, and a stopping rule chooses a stopping time; the chance of winning is
the proportion at the moment you stop, whose expectation is the initial proportion by the optional stopping
theorem (the deck is finite, so the conditions hold). A backward induction over all counts of red and black
remaining confirms it, and so does a simulation of a rule that waits for black cards to run ahead.
\iqlookfor{recognising a martingale, and not being seduced by rules that ``wait for a good moment''.}
\end{solution}

\begin{solution}{iq:iv:the-final-round-and-trading-games:10}
First candidate: ``Quoted 17 at 23 throughout; centre 1 below fair; never updated on cards or trades.''
Second: ``Tightened after each trade; sold three times to the same informed player; did not infer from flow.''
Third: ``Lifted the offer after two high cards were shown, citing the removal; tracked position; price near
fair.'' Pass the third. The first two may be strong in other interviews, but the game measures updating and
awareness of informed flow, and on those dimensions the notes are specific and negative; a note that quotes
numbers and behaviour is what a committee can use.
\iqlookfor{concrete, behavioural notes tied to the rubric, and a decision that follows from them.}
\end{solution}
